\section*{2026 年 3 月 16 日作业}
\phantomsection
\addcontentsline{toc}{section}{2026 年 3 月 16 日作业}
\noindent\textbf{截止时间：2026 年 3 月 23 日 13:30}
\setcounter{problem}{0}

\begin{problem}
    若 $\K$ 是 $\C$ 的子域，且 $\R$ 是 $\K$ 的子域，证明：$\K=\R$ 或者 $\K=\C$。

    注：本题说明不存在介于实数域和复数域之间的域。

    \begin{solution}
        设 $\R\subseteq \K\subseteq \C$。

        若 $\K=\R$，结论成立。

        下面设 $\K\neq \R$。则存在 $z\in \K\setminus \R$。由于 $z\in\C$，可写成
        \[
            z=a+b\ii,\qquad a,b\in\R.
        \]
        又因为 $z\notin\R$，所以 $b\neq 0$。

        由 $\R\subseteq \K$，知 $a,b\in \K$。于是
        \[
            z-a=b\ii\in \K.
        \]
        又因 $b\neq 0$ 且 $b\in \K$，故 $b^{-1}\in \K$。所以
        \[
            \ii=b^{-1}(z-a)\in \K.
        \]

        现任取任意 $w\in\C$。则存在 $u,v\in\R$ 使
        \[
            w=u+v\ii.
        \]
        由于 $\R\subseteq \K$ 且 $\ii\in \K$，故 $u,v,\ii\in \K$，从而
        \[
            w=u+v\ii\in \K.
        \]
        因此 $\C\subseteq \K$。

        另一方面，本来就有 $\K\subseteq \C$，所以
        \[
            \K=\C.
        \]

        综上，
        \[
            \K=\R\quad\text{或}\quad \K=\C.
        \]
        证毕。
    \end{solution}
\end{problem}

\begin{problem}
    分母有理化：
    \[
        \alpha=\frac{1}{2+3\sqrt[3]{3}+\sqrt[3]{9}},
        \qquad
        \beta=\frac{1}{1+2\sqrt2+2\sqrt3-3\sqrt6}.
    \]

    \begin{solution}
        先求 $\alpha$。设
        \[
            \omega=\sqrt[3]{3},\qquad \omega^2=\sqrt[3]{9},\qquad \omega^3=3.
        \]
        在基 $\left(1,\omega,\omega^2\right)$ 下，“乘以 $\omega$” 的线性变换矩阵为
        \[
            A=
            \begin{bmatrix}
                0 & 0 & 3\\
                1 & 0 & 0\\
                0 & 1 & 0
            \end{bmatrix},
            \qquad
            A^2=
            \begin{bmatrix}
                0 & 3 & 0\\
                0 & 0 & 3\\
                1 & 0 & 0
            \end{bmatrix}.
        \]
        因此 $2+3\omega+\omega^2$ 对应的矩阵为
        \[
            M=2I+3A+A^2=
            \begin{bmatrix}
                2 & 3 & 9\\
                3 & 2 & 3\\
                1 & 3 & 2
            \end{bmatrix}.
        \]
        于是
        \[
            \alpha=\frac1{2+3\omega+\omega^2}
        \]
        对应于 $M^{-1}$。计算得
        \[
            M^{-1}=
            \begin{bmatrix}
                -\frac5{44} & \frac{21}{44} & -\frac9{44}\\[2mm]
                -\frac3{44} & -\frac5{44} & \frac{21}{44}\\[2mm]
                \frac7{44} & -\frac3{44} & -\frac5{44}
            \end{bmatrix}.
        \]

        而一般地，
        \[
            p+q\omega+r\omega^2
            \longmapsto
            \begin{bmatrix}
                p & 3r & 3q\\
                q & p & 3r\\
                r & q & p
            \end{bmatrix},
        \]
        故由 $M^{-1}$ 的第一列可读出
        \[
            p=-\frac5{44},\qquad q=-\frac3{44},\qquad r=\frac7{44}.
        \]
        所以
        \[
            \boxed{
                \alpha
                =-\frac5{44}-\frac3{44}\sqrt[3]{3}+\frac7{44}\sqrt[3]{9}
                =\frac{-5-3\sqrt[3]{3}+7\sqrt[3]{9}}{44}
            }.
        \]

        再求 $\beta$。设
        \[
            u=\sqrt2,\qquad v=\sqrt3,\qquad uv=\sqrt6.
        \]
        在基 $\left(1,u,v,uv\right)$ 下，“乘以 $u$” 与 “乘以 $v$” 的矩阵分别为
        \[
            U=
            \begin{bmatrix}
                0 & 2 & 0 & 0\\
                1 & 0 & 0 & 0\\
                0 & 0 & 0 & 2\\
                0 & 0 & 1 & 0
            \end{bmatrix},
            \qquad
            V=
            \begin{bmatrix}
                0 & 0 & 3 & 0\\
                0 & 0 & 0 & 3\\
                1 & 0 & 0 & 0\\
                0 & 1 & 0 & 0
            \end{bmatrix}.
        \]
        因此 $1+2u+2v-3uv$ 对应的矩阵为
        \[
            N=I+2U+2V-3UV=
            \begin{bmatrix}
                1 & 4 & 6 & -18\\
                2 & 1 & -9 & 6\\
                2 & -6 & 1 & 4\\
                -3 & 2 & 2 & 1
            \end{bmatrix}.
        \]
        于是
        \[
            \beta=\frac1{1+2u+2v-3uv}
        \]
        对应于 $N^{-1}$。计算得
        \[
            N^{-1}=
            \begin{bmatrix}
                -\frac{31}{7} & -\frac{44}{7} & -\frac{54}{7} & -\frac{78}{7}\\[2mm]
                -\frac{22}{7} & -\frac{31}{7} & -\frac{39}{7} & -\frac{54}{7}\\[2mm]
                -\frac{18}{7} & -\frac{26}{7} & -\frac{31}{7} & -\frac{44}{7}\\[2mm]
                -\frac{13}{7} & -\frac{18}{7} & -\frac{22}{7} & -\frac{31}{7}
            \end{bmatrix}.
        \]

        而一般地，元素 $a+bu+cv+duv$ 在这组基下的坐标就是其“乘法矩阵”的第一列，因此由 $N^{-1}$ 的第一列可直接读出
        \[
            a=-\frac{31}{7},\qquad
            b=-\frac{22}{7},\qquad
            c=-\frac{18}{7},\qquad
            d=-\frac{13}{7}.
        \]
        故
        \[
            \boxed{
                \beta
                =-\frac{31}{7}-\frac{22}{7}\sqrt2-\frac{18}{7}\sqrt3-\frac{13}{7}\sqrt6
                =-\frac{31+22\sqrt2+18\sqrt3+13\sqrt6}{7}
            }.
        \]

        因此分母有理化结果为
        \[
            \boxed{
                \frac{1}{2+3\sqrt[3]{3}+\sqrt[3]{9}}
                =\frac{-5-3\sqrt[3]{3}+7\sqrt[3]{9}}{44}
            }
        \]
        和
        \[
            \boxed{
                \frac{1}{1+2\sqrt2+2\sqrt3-3\sqrt6}
                =-\frac{31+22\sqrt2+18\sqrt3+13\sqrt6}{7}
            }.
        \]
    \end{solution}
\end{problem}

\begin{problem}
    若 $K$ 是 $\C$ 的子域，在 $K^2$ 上定义运算
    \[
        \begin{bmatrix}
            a_1\\
            b_1
        \end{bmatrix}
        \boxplus
        \begin{bmatrix}
            a_2\\
            b_2
        \end{bmatrix}
        =
        \begin{bmatrix}
            a_1+a_2\\
            b_1+b_2+a_1a_2
        \end{bmatrix},
        \qquad
        k\boxtimes
        \begin{bmatrix}
            a\\
            b
        \end{bmatrix}
        =
        \begin{bmatrix}
            ka\\
            kb+k(k-1)a^2/2
        \end{bmatrix},
        \quad k\in K.
    \]
    证明：若把 $\boxplus$ 和 $\boxtimes$ 分别视为加法和数乘运算，那么 $K^2$ 在这两个运算下构成向量空间。

    \begin{solution}
        由于 $K\subseteq \C$，故 $\operatorname{char}K=0$，从而 $\frac12\in K$。

        定义映射
        \[
            \Phi:K^2\to K^2,\qquad
            \Phi\begin{pmatrix} a\\ b \end{pmatrix}
            =
            \begin{pmatrix}
                a\\
                b-\dfrac{a^2}{2}
            \end{pmatrix}.
        \]
        它的逆映射为
        \[
            \Phi^{-1}\begin{pmatrix} x\\ y \end{pmatrix}
            =
            \begin{pmatrix}
                x\\
                y+\dfrac{x^2}{2}
            \end{pmatrix},
        \]
        因而 $\Phi$ 是一个双射。

        设
        \[
            u=\begin{pmatrix} a_1\\ b_1 \end{pmatrix},\qquad
            v=\begin{pmatrix} a_2\\ b_2 \end{pmatrix}\in K^2,\qquad
            k\in K.
        \]
        由定义直接计算得
        \[
            \Phi(u\boxplus v)
            =
            \begin{pmatrix}
                a_1+a_2\\
                b_1+b_2+a_1a_2-\dfrac{(a_1+a_2)^2}{2}
            \end{pmatrix}
            =
            \begin{pmatrix}
                a_1+a_2\\
                \left(b_1-\dfrac{a_1^2}{2}\right)+\left(b_2-\dfrac{a_2^2}{2}\right)
            \end{pmatrix}
            =\Phi(u)+\Phi(v).
        \]
        同理，
        \[
            \Phi(k\boxtimes u)
            =
            \begin{pmatrix}
                ka_1\\
                kb_1+\dfrac{k(k-1)a_1^2}{2}-\dfrac{k^2a_1^2}{2}
            \end{pmatrix}
            =
            \begin{pmatrix}
                ka_1\\
                k\left(b_1-\dfrac{a_1^2}{2}\right)
            \end{pmatrix}
            =k\,\Phi(u).
        \]

        也就是说，题目中的运算满足
        \[
            u\boxplus v=\Phi^{-1}\bigl(\Phi(u)+\Phi(v)\bigr),\qquad
            k\boxtimes u=\Phi^{-1}\bigl(k\,\Phi(u)\bigr),
        \]
        其中右边的加法与数乘都是通常意义下 $K^2$ 的运算。

        因为通常的 $K^2$ 本来就是一个 $K$-向量空间，而上述两式说明 $(K^2,\boxplus,\boxtimes)$ 只是把这个向量空间经由双射 $\Phi$ 映到原集合上，所以 $(K^2,\boxplus,\boxtimes)$ 也是一个 $K$-向量空间。
    \end{solution}
\end{problem}

\begin{problem}
    若 $\K$ 是 $\C$ 的子域，证明：$\K^n$ 不能表示成其有限真子空间的并。
\end{problem}

\begin{problem}
    证明
    \[
        \left\{
        f(\mu)\in\mathbb{Q}[\mu]:
        (\mu-1)f(\mu+1)-(\mu+2)f(\mu)=0
        \right\}
    \]
    是 $\mathbb{Q}$ 上的向量空间，并求它的维数和一组基。
    \begin{solution}
        \begin{solnote}
            在这道题里，我们把 $(\mu-1)f(\mu+1)-(\mu+2)f(\mu)=0$ 写成分式： $\dfrac{f(\mu+1)}{f(\mu)} = \dfrac{\mu + 2}{\mu - 1}$，
            就能得到
            \begin{align}
                \dfrac{f(\mu + 4)}{f(\mu)}
                &=  \dfrac{f(\mu + 4)}{f(\mu + 3)}
                    \dfrac{f(\mu + 3)}{f(\mu + 2)}
                    \dfrac{f(\mu + 2)}{f(\mu + 1)}
                    \dfrac{f(\mu + 1)}{f(\mu)} \\
                &=  \dfrac{\mu + 5}{\mu + 2}
                    \dfrac{\mu + 4}{\mu + 1}
                    \dfrac{\mu + 3}{\mu}
                    \dfrac{\mu + 2}{\mu - 1} \\
                &= \dfrac{(\mu + 3)(\mu + 4)(\mu + 5)}{(\mu - 1)\mu(\mu + 1)}.\label{eq:mu-shift-ratio}
            \end{align}
            注意到令 $g(\mu) = (\mu - 1)\mu(\mu + 1)$ 的话，就有 $\dfrac{f(\mu + 4)}{f(\mu)} = \dfrac{g(\mu + 4)}{g(\mu)}$。
            
            所以我们至少可以猜测 $f(x)$ 应该与 $g(x)$ 有某种相似的形式。  
        \end{solnote}

        设 $V = \left\{
        f(\mu)\in\mathbb{Q}[\mu]:
        (\mu-1)f(\mu+1)-(\mu+2)f(\mu)=0
        \right\}$ ，我们先证明 $V$ 确实是 $\mathbb{Q}$ 上的向量空间。

        因为 $V\subseteq \mathbb{Q}[\mu]$，只需证明 $V$ 对线性组合封闭。任取
        \[
            f,g\in V,\qquad a,b\in\mathbb{Q},
        \]
        并令
        \[
            h(\mu)=af(\mu)+bg(\mu).
        \]
        由于 $f,g\in V$，有
        \[
            (\mu-1)f(\mu+1)-(\mu+2)f(\mu)=0,
        \]
        以及
        \[
            (\mu-1)g(\mu+1)-(\mu+2)g(\mu)=0.
        \]
        因而
        \begin{align*}
            (\mu-1)h(\mu+1)-(\mu+2)h(\mu)
            &= (\mu-1)\bigl(af(\mu+1)+bg(\mu+1)\bigr) \\
            &\quad - (\mu+2)\bigl(af(\mu)+bg(\mu)\bigr) \\
            &= a\bigl[(\mu-1)f(\mu+1)-(\mu+2)f(\mu)\bigr] \\
            &\quad + b\bigl[(\mu-1)g(\mu+1)-(\mu+2)g(\mu)\bigr] \\
            &= 0.
        \end{align*}
        所以 $h\in V$，故 $V$ 对线性组合封闭。
        因而 $V$ 是 $\mathbb{Q}[\mu]$ 的子空间，从而是 $\mathbb{Q}$ 上的向量空间。

        下面我们来求 $V$ 的维数和一组基。

        首先由 \eqref{eq:mu-shift-ratio} 式得，
        $\dfrac{f(\mu + 4)}{f(\mu)} = \dfrac{g(\mu + 4)}{g(\mu)}$。
        
        所以我们令 $h(\mu) = \dfrac{f(\mu)}{g(\mu)}$的话，就有 $h(\mu + 4) = h(\mu)$。
        
        又因为 $f(x)$ 和 $g(x)$ 都是多项式，所以 $h(x)$ 是周期性的有理函数。

        \begin{sollemma*}
            设 \(r(x)\in \mathbb{Q}(x)\) 是一个有理函数。若存在非零常数 \(T\in \mathbb{Q}\)，使得
            \[
                r(x+T)=r(x)
            \]
            对一切使两边都有定义的 \(x\) 成立，则 \(r(x)\) 必为常数。
        \end{sollemma*}

        \begin{proof}
            设
            \[
                r(x)=\frac{p(x)}{q(x)},
                \qquad p(x),q(x)\in\mathbb{Q}[x],\ q(x)\neq 0.
            \]
            记
            \[
                m=\deg p,\qquad n=\deg q.
            \]

            由于 \(r(x)\) 是周期函数，存在非零常数 \(T\) 使得
            \[
                r(x+T)=r(x)
            \]
            恒成立。下分三种情形讨论。

            \medskip

            \noindent\textbf{情形 1: \(m<n\).}

            此时
            \[
                \lim_{x\to\infty} r(x)=0.
            \]
            任取 \(x\) 使得 \(r(x)\) 有定义，由周期性可得对任意正整数 \(k\),
            \[
                r(x)=r(x+kT).
            \]
            当 \(k\to\infty\) 时，\(x+kT\to\infty\)，故
            \[
                r(x)=\lim_{k\to\infty}r(x+kT)=0.
            \]
            于是 \(r(x)\equiv 0\)，从而 \(r(x)\) 为常数。

            \medskip

            \noindent\textbf{情形 2: \(m=n\).}

            设 \(p(x)\) 与 \(q(x)\) 的首项系数分别为 \(a_m\) 与 \(b_n\)，则
            \[
                \lim_{x\to\infty} r(x)=\frac{a_m}{b_n}.
            \]
            同样地，任取 \(x\) 使得 \(r(x)\) 有定义，由周期性有
            \[
                r(x)=r(x+kT)\qquad (\forall k\in\mathbb{N}).
            \]
            令 \(k\to\infty\)，得
            \[
                r(x)=\lim_{k\to\infty}r(x+kT)=\frac{a_m}{b_n}.
            \]
            因此 \(r(x)\) 恒等于常数 \(\dfrac{a_m}{b_n}\)。

            \medskip

            \noindent\textbf{情形 3: \(m>n\).}

            此时当 \(x\to\infty\) 时，
            \[
                |r(x)|\to\infty.
            \]
            任取 \(x\) 使得 \(r(x)\) 有定义，由周期性仍有
            \[
                r(x)=r(x+kT)\qquad (\forall k\in\mathbb{N}).
            \]
            当 \(k\to\infty\) 时，\(x+kT\to\infty\)，于是
            \[
                |r(x+kT)|\to\infty.
            \]
            但左边又恒等于固定值 \(|r(x)|\)，矛盾。

            故情形 3 不可能发生。

            \medskip

            综上，只有前两种情形可能出现，而这两种情形都推出 \(r(x)\) 为常数。因此，周期有理函数必为常数。
        \end{proof}

        所以 $h(\mu)$ 必为常数，即 $\dfrac{f(\mu)}{g(\mu)} = c \Rightarrow f(\mu) = c(\mu -1)\mu(\mu + 1), \quad c \in \mathbb{Q}$。
        
        所以 $V$ 的维数为1,一个基为 $(\mu -1)\mu(\mu + 1)$。 
    \end{solution}
\end{problem}

\begin{problem}[选做]
    在 $\R^2$ 上定义乘法 $\times$：
    \[
        \begin{bmatrix}
            \alpha_1\\
            \beta_1
        \end{bmatrix}
        \times
        \begin{bmatrix}
            \alpha_2\\
            \beta_2
        \end{bmatrix}
        =
        \begin{bmatrix}
            \alpha_1\alpha_2-\beta_1\beta_2\\
            \alpha_1\beta_2+\alpha_2\beta_1
        \end{bmatrix}.
    \]
    则可构成一个与 $\C$ 同构的域，且
    \[
        \left\{
        \begin{bmatrix}
            \alpha\\
            0
        \end{bmatrix}
        :\alpha\in\R
        \right\}
    \]
    同构于 $\R$。证明：在 $\R^3$ 上不论如何定义乘法，都无法构成一个域使得
    \[
        \left\{
        \begin{bmatrix}
            \alpha\\
            0\\
            0
        \end{bmatrix}
        :\alpha\in\R
        \right\}
    \]
    同构于 $\R$。
\end{problem}

\begin{problem}[选做]
    若 $n$ 是正整数，$A\in\R^{n\times n}$ 是给定的矩阵，且对任何 $x,y\in\R^n$ 都有
    \[
        (y^\top Ax)^2=(x^\top Ay)^2.
    \]
    证明：$A$ 是对称矩阵或反对称矩阵。
    \begin{solution}
        因为 $y^TAx$ 和 $x^TA y$ 都是实数，所以有 $y^TAx = x^TAy$ 或 $y^TAx = -x^TAy$。
        
        又因为
        \[
            y^T Ax = (Ax)^Ty = x^T A^T y,
        \]
        所以有 $x^T(A^T \pm A)y = 0$ 对于任意 $x,y\in\R^n$ 都成立，

        所以 $A^T \pm A = 0$，即 $A$ 是对称矩阵或反对称矩阵。
    \end{solution}
\end{problem}
